\section{Problem}
\label{section:12-problem}

\gai covers a broad collection of tasks with varying computational requirements. 
Most organizations that employ \gai will need to solve problems using tasks from across the \gai taxonomy. 
Selecting the optimal system configuration, hardware accelerators and implementations of \gai tasks is essential to effectively addressing the breadth of real-world problems, particularly when faced with time, cost and other constraints. 
Organizations may be required to trade-off between the effectiveness, efficiency and cost of solutions, and need to understand what the implications of those tradeoffs are. 

Competitive benchmarks and workload analyses offer two paths to understand these trade-offs, but either provide too little information or require invasive instrumentation to characterize the System Under Test (SUT).

This thesis proposes a hybrid solution to the problem of \gai benchmarking, where we draw from both existing benchmarking and workload analysis studies to develop an approach that allows for an investigator to trade-off between effectiveness and efficiency when testing a system.
% A hybrid benchmark must provide the dataset coverage to test variance from differing input shapes and sizes, support the varying abstractions and representations that drive implementations and collect the metrics required to make informed decisions about the tradeoffs between effectiveness and efficiency. 
% It must allow one component to vary while the others remain fixed, and decouple the components sufficiently so that any desired combination of components can be tested. 
We develop a benchmarking framework for \gai to bridge the gap between competitive benchmarks and invasive workload analyses that limits access to the Input/Output interface and maximizes the freedom of a user to modify the system under test.
Our approach is problem-centric, system-agnostic and gives the user a mechanism for understanding of how system and implementation changes impact the efficiency, effectiveness and cost of their system. 
The design draws on the empirical study of algorithms, constraint theory and the significant body of knowledge in benchmarking and analyzing workloads of \gai to develop a formal methodology for developing \gai benchmark problems. 


% Selecting an appropriate \gai implementation and runtime environment for a given task and dataset requires an understanding of the efficiency and effectiveness of particular implementations in different environments under different inputs.
% To further complicate the decision, multiple \gai tasks may need to coexist in one environment to accomplish a broader goal, which makes the tradeoffs even more challenging to balance.
% Users of \gai may need to achieve multiple \gai tasks in the same environment, so they will need to understand the tradeoffs in efficiency and effectiveness of completing each task for given combinations of implementations and environments while staying within some cost constraint.



% \nrscomment{you already introduced benchmarks above. this chunk is a bunch of short sentences that are quite redundant.i think youre better off replacing all of this with a succinct summary of why its so difficult to measure the right thing to make the tradeoff appropriately (i.e. why current benchmarks are insufficient). you were building towards that with the discussion of multiple tasks in one environment, and how the tradeoff becomes more complicated in that case, then just kind of left it on the ground. i assume the coexisting competing tasks thing is super relevant for things like GNNs and other hybrid problems-- can you take that further and really convey why this is so hard that no one has done it correctly before? i thinky ou need a few sentences here to drive home the difficulty of the challenge at hand, and then bridge that to the next paragraph that proposes your solution to it.}
% Benchmarking is a process by which systems or components are compared according to some criteria, that provide a point of reference to measure the relative performance~\cite{Kistowski2015}.
% It is used in many fields and for many different tasks, but here specifically refers to evaluating and comparing \textit{effectiveness} and \textit{efficiency}, (which will be collectively called \textit{performance}) and \textit{cost}.
% Effectiveness is a measure of `how well' a problem is solved and typically uses metrics like accuracy, precision, recall, and error rate.
% Ideally, a system should maximize its effectiveness. 
% Efficiency is a measure of `how quickly' a system solves a problem and uses metrics reflecting execution time, throughput and resource usage statistics. 
% Ideally, a system should maximize its efficiency (minimize its cost).
% Cost is a measure of `how expensive' it is for a system to solve a problem and can use purely financial metrics, or proxy estimates like power consumption.
% Ideally, a system should minimize cost. 
% The problem is then, what should a comprehensive \gai benchmark to effectively compare the effectiveness, efficiency and cost of a system look like?

% \osullikomment{Diagram here about the tension between effectiveness, efficiency and cost}



% \nrscomment{@nsch look at this paragraph again once the above one is fleshed out. need to more concisely and clearly convey the contribution of the thesis.}



% \subsection{Why Benchmark}
% \nrscomment{this also feels like problem/motivation content. I didn't read it in detail, but from the high level look at what its saying, id suggest reframing it slightly to get away from the user/developer focus and instead attach it as "the types of things a benchmark should enable are x, y decision making, which requires z information. That can be supported with examples of how someone doing procurement might use that info, but i think that can be secondary to the motivation youre trying to accomplish, rather than the primary focus of this chunk. If you do use an example to demonstrate the utility and can make it more specific to intelligence context that seven better. this whole thing is starting to stray from the IandS perspective, so i think sprinkling it in whenever you need to use examples will help keep it cohesive.}

% Benchmarking has long proven its utility in assisting users to make decisions about which systems to procure or frameworks to use, and helping developers understand how to optimize their systems, abstractions and algorithms in future iterations. 
% Here, a \textit{user} represents somebody deciding which system they should procure to meet their business need. 
% Users can also be software engineers who are making decisions about which framework, abstraction, programming model or algorithm to implement in their system. 
% In either case, a user wishes to measure the relative performance of candidate components to make an informed decision about the \gai components they are integrating it into a larger system to address a business problem~\cite{Carpinetti2002}.
% A \textit{developer} represents someone working to design and build the fundamental hardware, programs, algorithms and representations that comprise the \gai stack. 
% Developers are focused on building and optimizing the individual components of \gai systems, and benchmarks provide a starting point to target more detailed workload analyses to identify bottlenecks and inform future development~\cite{Saavedra1996, Hooker1994, Hooker1995}. 

% Whether for a user or a developer, that benchmarking process relies on measuring the \textit{effectiveness} and \textit{efficiency} of systems to reduce their uncertainty about the aspects of \gai relevant to them.
% Benchmarks for high-performance computing like the Standard Performance Evaluation Corporation (SPEC) benchmarks~\cite{Hennessy2019}, for solving linear equations like the LINPACK Top500~\cite{Dongarra2003}, for satisfiability problems like SATLIB~\cite{Hoos2000}, and far more have been essential in driving acquisition decisions for users and guiding optimization efforts for developers across all these fields. 
% Benchmarks are a useful tool, and section~\ref{section:32-related-benchmarking-gai} shows how they have been useful for \gai specifically.



